% Meta-Diffusion technical report aligned with the executable experiment plan.
% Compile from docs/ with:
%   pdflatex -interaction=nonstopmode -output-directory=../output/pdf iclr_executable_plan_results_2026-09-21.tex
\documentclass{article}

\usepackage{iclr2026_conference,times}
\usepackage{amsmath,amssymb}
\usepackage{booktabs}
\usepackage{graphicx}
\usepackage{xcolor}
\usepackage{enumitem}
\usepackage{etoolbox}
\usepackage{placeins}
\usepackage{hyperref}
\usepackage{url}

\newcommand{\KT}{K_T}
\newcommand{\MS}{M_S}
\newcommand{\SW}{\mathrm{SW}}
\newcommand{\MMD}{\mathrm{MMD}}
\newcommand{\CI}{95\%\ \mathrm{CI}}
\newcommand{\sigstar}{^{*}}
\newcommand{\method}{\textsc{Meta-Diffusion}}
\setlist{nosep,leftmargin=1.35em}
\setlength{\emergencystretch}{3em}
\AfterEndEnvironment{table}{\FloatBarrier}

\title{Meta-Diffusion:\\
Plan-Aligned Results\\
and Fitzpatrick17k Source Specificity}
\author{Anonymous Authors}

\begin{document}
\maketitle

\begin{abstract}
We report the completion of a pre-registered experiment sequence designed to distinguish
reusable population transport from genuinely source-task-specific transport.  The sequence
audits the data and frozen method (P0), raises Fitzpatrick17k task resolution from diseases to
pathology families (P1), introduces hierarchical wrong-source controls (P2), tests a
source-only task-separability explanation (P3), repeats the family protocol under a stronger
Fitzpatrick I--III to V--VI shift (P4), audits the completed CIFAR package (P5), and applies a
strict gate to downstream demographic augmentation (P6).  On Fitzpatrick17k, correct-source
transport improves energy MMD over a population-only control, but is statistically
indistinguishable from wrong-task transport under both phenotype definitions.  The frozen
family classifier also fails its pre-specified per-family validity rule, so semantic nulls are
classified as inconclusive rather than equivalence.  Source-only task separability does not
provide a supported explanation.  In contrast, audited CIFAR evidence retains a clear
correct-versus-wrong-source advantage at $\KT=0$, survives FiLM conditioning, and supports an
approximately linear task-space relation.  The resulting claim is therefore deliberately
narrow: source-specific coordinate transport is supported in the controlled CIFAR setting,
whereas the frozen method learns primarily a shared phenotype offset on Fitzpatrick17k.
\end{abstract}

\section{Scope and Reporting Convention}

The core model, training recipe, linear transport, and no-refinement headline inference were
frozen before the new family-level outcomes were inspected.  No new head, larger coordinate,
nonlinear transport, or architecture rescue was introduced.  All within-task comparisons use
the same support draw, real target query, and DDIM noise.  Repeats reduce Monte Carlo noise but
are not treated as independent semantic tasks.

Raw-performance tables report absolute means $\pm$ 95\% confidence-interval half-widths.  A
separate comparison table reports absolute paired differences.  For lower-is-better
distances, the comparison is
\begin{equation}
  \Delta_d=d(\text{comparator},\mathcal D_T)
          -d(\text{correct source+transport},\mathcal D_T),
\end{equation}
so positive values favour correct-source transport.  For higher-is-better semantic or
phenotype consistency, the subtraction is reversed, again making positive values favour
correct transport.  A superscript $\sigstar$ is attached directly to a paired result only
when its 95\% interval excludes zero.  No star is placed on an unpaired raw mean.

The reported measurements are: phenotype target share and semantic consistency
(higher is better), and fixed phototype-statistic distance, sliced Wasserstein distance
(SW), and energy MMD (lower is better).  A frozen predictor is a measurement instrument, not
ground truth.  Failure of its independent validity gate prevents a semantic claim and is not
replaced post hoc by a distributional metric.

\section{P0: Prerequisite and Reproducibility Gate}

The experiment used the verified SFU Fitzpatrick17k-Renamed archive as a byte source.  Images
were matched to the official CSV using the supplied MD5 hashes and filtered with the
Fitzpatrick17k-C allow-list.  Known near-duplicate components were collapsed before task
allocation.  The official schema contains no subject identifier; consequently patient-level
leakage cannot be checked directly, and hash/near-duplicate filtering is the strongest
available safeguard.

\begin{table}[h]
\caption{P0 dataset, environment, and frozen-protocol audit.}
\label{tab:p0}
\centering
\small
\begin{tabular}{lr}
\toprule
Audit item & Verified value \\
\midrule
Official CSV / archive rows & 16,577 / 16,577 \\
Byte-verified rows retained by Fitzpatrick17k-C & 11,394 \\
Retained rows with Fitzpatrick I--VI & 11,098 \\
P1 eligible pathology families (target IV--VI) & 6 \\
P4 eligible pathology families (target V--VI) & 5 \\
P1 task-level folds / P4 cyclic folds & 3 / 5 \\
Training / evaluation seeds & 12,345 / 4,321 \\
Python / PyTorch / CUDA & 3.12.13 / 2.11.0 / 12.8 \\
GPU & NVIDIA RTX 5080, 16,303 MiB \\
Frozen coordinate dimension / source budget & $k=16$ / $\MS=32$ \\
\bottomrule
\end{tabular}
\end{table}

All manifests passed on-disk, population-side, taxonomy, duplicate, and capacity checks.
Target support is disjoint from target query and nested across $\KT\in\{1,2,5,10,20\}$.
The family-level semantic-instrument check remains an outcome-validity gate and was therefore
not used to select tasks or folds.

\paragraph{P0 conclusion.}
The prerequisite gate passed.  Data provenance, exclusions, task folds, frozen checkpoints,
seeds, comparison direction, and independent statistical units were fixed before the new
outcomes were inspected.

\section{P1: Pathology-Family Semantic Resolution}

The task is the official \texttt{nine\_partition\_label}; source is Fitzpatrick I--III and
target is IV--VI.  Six families meet the pre-model requirement of at least 64 clean source
and 45 clean target images.  Three task-level folds test every family exactly once.  Model
selection uses only the validation family in each fold and chooses training steps 16,000,
10,000, and 10,000.  Evaluation uses 96 generated images per cell, DDIM-50, eight paired
repeats, and no target-time refinement.

\subsection{Absolute zero-shot performance}

\begin{table}[h]
\caption{P1 raw performance at $\KT=0$ over six outer-test pathology families.  Entries are
absolute family means $\pm$ 95\% CI.  Phenotype share and family consistency are
higher-is-better; all three distances are lower-is-better.}
\label{tab:p1raw}
\centering
\scriptsize
\resizebox{\textwidth}{!}{%
\begin{tabular}{lccccc}
\toprule
Method & Phenotype share $\uparrow$ & Family consistency $\uparrow$ & Statistic dist. $\downarrow$ & SW $\downarrow$ & Energy MMD $\downarrow$ \\
\midrule
Source reuse & $0.0000\pm0.0000$ & $0.1651\pm0.3227$ & $1.7452\pm0.2089$ & $0.3081\pm0.0192$ & $5.7946\pm1.0078$ \\
Population only & $0.0000\pm0.0000$ & $0.1641\pm0.3216$ & $3.0504\pm0.2209$ & $0.2683\pm0.0110$ & $6.8472\pm1.0089$ \\
Wrong-family source $+$ relation & $0.0000\pm0.0000$ & $0.1656\pm0.3230$ & $1.7185\pm0.2170$ & $0.3064\pm0.0200$ & $5.7468\pm1.0114$ \\
Correct-family source $+$ transport & $0.0002\pm0.0004$ & $0.1656\pm0.3230$ & $1.7188\pm0.2178$ & $0.3063\pm0.0199$ & $5.7436\pm1.0026$ \\
\bottomrule
\end{tabular}}
\end{table}

\begin{table}[h]
\caption{P1 paired effects at $\KT=0$.  Every entry is oriented so positive values favour
correct-family transport.  A star means the paired 95\% CI excludes zero.}
\label{tab:p1paired}
\centering
\scriptsize
\resizebox{\textwidth}{!}{%
\begin{tabular}{lccccc}
\toprule
Comparator ($\Delta>0$: correct-family transport) & Phenotype share & Family consistency & Statistic dist. & SW & Energy MMD \\
\midrule
Source reuse & $+0.00022\pm0.00043$ & $+0.00043\pm0.00054$ & $+0.02635\pm0.01482\sigstar$ & $+0.00178\pm0.00218$ & $+0.05099\pm0.05141$ \\
Population only & $+0.00022\pm0.00043$ & $+0.00152\pm0.00191$ & $+1.33162\pm0.08916\sigstar$ & $-0.03795\pm0.02043\sigstar$ & $+1.10360\pm0.49696\sigstar$ \\
Wrong-family source $+$ relation & $+0.00022\pm0.00043$ & $+0.00000\pm0.00000$ & $-0.00029\pm0.00160$ & $+0.00007\pm0.00091$ & $+0.00317\pm0.01925$ \\
\bottomrule
\end{tabular}}
\end{table}

The semantic evaluator obtains overall fold accuracies of 67.6--68.7\%, but its minimum
outer-test-family accuracies are 0.0\%, 0.0\%, and 15.3\%, below the locked 25\% per-family
threshold.  Hence the primary correct-minus-wrong semantic effect of exactly zero is
uninformative rather than evidence of equivalence.  The independently valid secondary
correct-versus-wrong effects all include zero.

\subsection{Target-only data-efficiency comparison}

Because the headline correct, wrong, source-reuse, and population-only coordinates receive no
target-time refinement, they are invariant to $\KT$.  Nonzero target budgets therefore test
the target-only comparator rather than further adapting transport.

\begin{table}[h]
\caption{P1 target-only raw distances and paired effects against the fixed correct-family
transport.  Within each metric, the first row is target-only raw performance and the second
is $\Delta=d(\text{target only})-d(\text{correct transport})$; positive supports correct
transport.}
\label{tab:p1budget}
\centering
\scriptsize
\resizebox{\textwidth}{!}{%
\begin{tabular}{llccccc}
\toprule
Metric & Quantity & $\KT=1$ & $\KT=2$ & $\KT=5$ & $\KT=10$ & $\KT=20$ \\
\midrule
Statistic dist. & Target only raw & $1.7322\pm0.2193$ & $1.7306\pm0.2201$ & $1.7412\pm0.2117$ & $1.7369\pm0.2147$ & $1.7356\pm0.2158$ \\
& Paired $\Delta$ ($+=$ correct) & $+0.0134\pm0.0227$ & $+0.0118\pm0.0127$ & $+0.0224\pm0.0119\sigstar$ & $+0.0180\pm0.0080\sigstar$ & $+0.0168\pm0.0092\sigstar$ \\
\midrule
SW & Target only raw & $0.3045\pm0.0194$ & $0.3055\pm0.0191$ & $0.3041\pm0.0190$ & $0.3047\pm0.0204$ & $0.3039\pm0.0194$ \\
& Paired $\Delta$ ($+=$ correct) & $-0.0018\pm0.0037$ & $-0.0007\pm0.0012$ & $-0.0022\pm0.0018\sigstar$ & $-0.0016\pm0.0014\sigstar$ & $-0.0024\pm0.0020\sigstar$ \\
\midrule
Energy MMD & Target only raw & $5.7208\pm0.9739$ & $5.7303\pm0.9850$ & $5.6749\pm0.9871$ & $5.7091\pm1.0350$ & $5.6900\pm1.0112$ \\
& Paired $\Delta$ ($+=$ correct) & $-0.0228\pm0.0745$ & $-0.0133\pm0.0315$ & $-0.0687\pm0.0379\sigstar$ & $-0.0345\pm0.0374$ & $-0.0536\pm0.0371\sigstar$ \\
\bottomrule
\end{tabular}}
\end{table}

The distance metrics disagree in direction, and target-only itself changes little with
$\KT$.  Consequently, no scientifically meaningful ``transport at $\KT=2$ equals target-only
at $\KT=10$'' statement is supported.

\paragraph{P1 conclusion.}
Raising semantic resolution to pathology families does not recover a source-specific
advantage.  The primary result is inconclusive because its evaluator fails the validity gate;
valid distributional diagnostics independently show no resolved correct-versus-wrong effect.

\section{P2: Hierarchical Wrong-Source Control}

The frozen disease-level checkpoint is evaluated on four test diseases with correct disease,
wrong disease in the same pathology family, wrong disease in a different family, source
reuse, and population-only conditions.  Wrong sources are deterministic and selected without
using model outcomes.

\begin{table}[h]
\caption{P2 raw zero-shot performance over four disease tasks.  Consistency measures are
higher-is-better; distances are lower-is-better.}
\label{tab:p2raw}
\centering
\scriptsize
\resizebox{\textwidth}{!}{%
\begin{tabular}{lccccc}
\toprule
Method & Phenotype share $\uparrow$ & Disease consistency $\uparrow$ & Statistic dist. $\downarrow$ & SW $\downarrow$ & Energy MMD $\downarrow$ \\
\midrule
Source reuse & $0.1390\pm0.0103$ & $0.0199\pm0.0224$ & $2.5248\pm0.4420$ & $0.3642\pm0.0448$ & $7.1765\pm1.8278$ \\
Population only & $0.0000\pm0.0000$ & $0.0000\pm0.0000$ & $4.0465\pm0.3111$ & $0.2867\pm0.0094$ & $8.4246\pm1.4825$ \\
Wrong disease, same family & $0.1423\pm0.0101$ & $0.0199\pm0.0224$ & $2.5020\pm0.4478$ & $0.3638\pm0.0454$ & $7.1466\pm1.8509$ \\
Wrong disease, different family & $0.1426\pm0.0098$ & $0.0205\pm0.0231$ & $2.5025\pm0.4462$ & $0.3639\pm0.0454$ & $7.1444\pm1.8488$ \\
Correct disease $+$ transport & $0.1419\pm0.0094$ & $0.0199\pm0.0221$ & $2.5030\pm0.4480$ & $0.3638\pm0.0456$ & $7.1429\pm1.8566$ \\
\bottomrule
\end{tabular}}
\end{table}

\begin{table}[h]
\caption{P2 hierarchical paired effects at $\KT=0$.  Positive values favour the first named
condition.  A star denotes a paired 95\% interval excluding zero.}
\label{tab:p2paired}
\centering
\scriptsize
\resizebox{\textwidth}{!}{%
\begin{tabular}{lccccc}
\toprule
Comparison ($\Delta>0$: first condition) & Phenotype share & Disease consistency & Statistic dist. & SW & Energy MMD \\
\midrule
Correct vs. same-family wrong & $-0.00033\pm0.00283$ & $+0.00000\pm0.00104$ & $-0.00101\pm0.00777$ & $+0.00009\pm0.00062$ & $+0.00365\pm0.01937$ \\
Same-family vs. different-family wrong & $-0.00033\pm0.00191$ & $-0.00065\pm0.00074$ & $+0.00053\pm0.00532$ & $+0.00004\pm0.00038$ & $-0.00219\pm0.01319$ \\
Correct vs. population only & $+0.14193\pm0.00938\sigstar$ & $+0.01986\pm0.02214$ & $+1.54346\pm0.14526\sigstar$ & $-0.07706\pm0.03919\sigstar$ & $+1.28164\pm0.52335\sigstar$ \\
\bottomrule
\end{tabular}}
\end{table}

The disease instrument reaches 20.8\% overall but only 53.1\%, 0\%, 6.2\%, and 0\% on the
four individual test diseases.  It is therefore unsuitable for a strong semantic ordering.
The correct/same-family/different-family distance contrasts are also numerically small and all
include zero.  Population-only again produces low pixel SW while failing phenotype and energy
MMD checks, demonstrating why no single pixel metric is sufficient.

\paragraph{P2 conclusion.}
Neither $\text{correct}>\text{same-family wrong}>\text{different-family wrong}$ nor a
family-only ordering is resolved.  The null hierarchy is not an equivalence result.

\section{P3: Outcome-Independent Source-Task Separability}

Official ImageNet ResNet-18 penultimate features were extracted from source-population
images only, L2-normalised per image, and averaged by disease.  Separation is Euclidean
distance to the pre-assigned different-family wrong source from P2.

\begin{table}[h]
\caption{P3 source-only separation and held-out semantic effect.}
\label{tab:p3}
\centering
\scriptsize
\begin{tabular}{lcc}
\toprule
Target disease & Separation & Correct-minus-wrong effect \\
\midrule
Basal cell carcinoma & 0.2488 & 0.0000 \\
Neurofibromatosis & 0.2807 & 0.0000 \\
Pityriasis rosea & 0.3289 & $-0.0026$ \\
Vitiligo & 0.2971 & 0.0000 \\
\bottomrule
\end{tabular}
\end{table}

Spearman $\rho=-0.775$ with a task-bootstrap interval $[-1.000,-0.775]$.  Only 6,793 of
10,000 bootstrap draws have defined rank variance because $n=4$ and three outcome effects are
exactly zero.  Moreover, the outcome is measured by the weak disease instrument described
above.

\paragraph{P3 conclusion.}
Greater source-task separability does not predict a larger correct-source advantage in the
locked analysis.  The result is unsupported/unresolved, not a negative causal mechanism, and
task identifiability should not be the paper's central explanation.

\section{P4: Stronger Phenotype Separation}

P4 was triggered because P1 was inconclusive and the pre-outcome counts were adequate.  The
only domain change is target Fitzpatrick V--VI; source remains I--III.  Five families are
tested in five cyclic outer folds, with each family's two appearances averaged before
inference.  Validation-only selection chooses steps 16,000, 4,000, 4,000, 4,000, and 10,000.
The architecture is a 32$\times$32 additive-basis small U-Net (0.72M parameters), set encoder
(0.33M), linear transport (0.3K), coordinate dimension 16, and no target-time refinement.

\begin{table}[h]
\caption{P4 raw performance at $\KT=0$ over five pathology families.  The stronger target is
Fitzpatrick V--VI.}
\label{tab:p4raw}
\centering
\scriptsize
\resizebox{\textwidth}{!}{%
\begin{tabular}{lccccc}
\toprule
Method & Phenotype share $\uparrow$ & Family consistency $\uparrow$ & Statistic dist. $\downarrow$ & SW $\downarrow$ & Energy MMD $\downarrow$ \\
\midrule
Source reuse & $0.0311\pm0.0359$ & $0.2033\pm0.3807$ & $1.8099\pm0.2613$ & $0.2649\pm0.0446$ & $4.5031\pm0.7991$ \\
Population only & $0.0000\pm0.0000$ & $0.1997\pm0.3909$ & $2.2577\pm0.4548$ & $0.2527\pm0.0275$ & $5.4122\pm1.1998$ \\
Wrong-task source $+$ relation & $0.0320\pm0.0369$ & $0.2033\pm0.3807$ & $1.7943\pm0.2765$ & $0.2637\pm0.0454$ & $4.4233\pm0.8218$ \\
Correct source $+$ transport & $0.0315\pm0.0363$ & $0.2031\pm0.3808$ & $1.7950\pm0.2769$ & $0.2638\pm0.0461$ & $4.4244\pm0.8362$ \\
\bottomrule
\end{tabular}}
\end{table}

\begin{table}[h]
\caption{P4 paired effects at $\KT=0$.  All entries are oriented so positive values favour
correct-source transport.}
\label{tab:p4paired}
\centering
\scriptsize
\resizebox{\textwidth}{!}{%
\begin{tabular}{lccccc}
\toprule
Comparator ($\Delta>0$: correct transport) & Phenotype share & Family consistency & Statistic dist. & SW & Energy MMD \\
\midrule
Source reuse & $+0.00039\pm0.00051$ & $-0.00013\pm0.00102$ & $+0.01494\pm0.01616$ & $+0.00117\pm0.00272$ & $+0.07874\pm0.08272$ \\
Population only & $+0.03151\pm0.03635$ & $+0.00339\pm0.01614$ & $+0.46273\pm0.46930$ & $-0.01112\pm0.02383$ & $+0.98779\pm0.70849\sigstar$ \\
Wrong-task source $+$ relation & $-0.00052\pm0.00074$ & $-0.00013\pm0.00026$ & $-0.00068\pm0.00273$ & $-0.00003\pm0.00134$ & $-0.00107\pm0.03907$ \\
\bottomrule
\end{tabular}}
\end{table}

The phenotype instrument is useful on outer-test families (81.2--92.7\%, chance 50\%).  The
semantic instrument has 69.9--70.5\% overall accuracy, yet worst-family accuracies of 0.0\%,
0.0\%, 0.0\%, 12.7\%, and 2.5\%, so every fold fails the locked per-family rule.  Pooled KID
is consequently skipped.  Correct transport significantly improves energy MMD over
population-only, but the correct-versus-wrong-task effect is $-0.0011\pm0.0391$ and is
unresolved.  The same near-zero pattern holds for every other valid secondary metric.

\begin{table}[h]
\caption{P4 target-only raw distances and paired effects against fixed correct transport.
For each metric, $\Delta=d(\text{target only})-d(\text{correct transport})$; positive values
support correct transport.}
\label{tab:p4budget}
\centering
\scriptsize
\resizebox{\textwidth}{!}{%
\begin{tabular}{llccccc}
\toprule
Metric & Quantity & $\KT=1$ & $\KT=2$ & $\KT=5$ & $\KT=10$ & $\KT=20$ \\
\midrule
Statistic dist. & Target only raw & $1.8082\pm0.2694$ & $1.8050\pm0.2706$ & $1.8044\pm0.2691$ & $1.7997\pm0.2702$ & $1.7972\pm0.2689$ \\
& Paired $\Delta$ ($+=$ correct) & $+0.0132\pm0.0109\sigstar$ & $+0.0100\pm0.0121$ & $+0.0094\pm0.0086\sigstar$ & $+0.0047\pm0.0071$ & $+0.0022\pm0.0106$ \\
\midrule
SW & Target only raw & $0.2597\pm0.0413$ & $0.2622\pm0.0440$ & $0.2624\pm0.0443$ & $0.2632\pm0.0449$ & $0.2630\pm0.0447$ \\
& Paired $\Delta$ ($+=$ correct) & $-0.0041\pm0.0052$ & $-0.0016\pm0.0022$ & $-0.0013\pm0.0020$ & $-0.0006\pm0.0014$ & $-0.0008\pm0.0017$ \\
\midrule
Energy MMD & Target only raw & $4.3571\pm0.7399$ & $4.3970\pm0.7715$ & $4.4157\pm0.7973$ & $4.4171\pm0.7974$ & $4.3974\pm0.7895$ \\
& Paired $\Delta$ ($+=$ correct) & $-0.0673\pm0.1207$ & $-0.0274\pm0.0679$ & $-0.0086\pm0.0460$ & $-0.0073\pm0.0406$ & $-0.0270\pm0.0576$ \\
\bottomrule
\end{tabular}}
\end{table}

\paragraph{P4 conclusion.}
A larger phenotype gap does not recover disease-family-specific source use.  The model learns
a reusable population shift relative to population-only on energy MMD, but obtains the same
shift from the wrong source task.  Under the registered stop rule, demographic model tuning
ends here.

\section{P5: Frozen CIFAR Package}

P5 was audited from existing raw per-episode artifacts rather than re-estimated
opportunistically.  The saved files contain matched episodes and paired intervals.

\begin{table}[h]
\caption{Audited CIFAR evidence.  All stated improvements are paired; starred rows have a
95\% interval excluding zero in the underlying artifact.}
\label{tab:p5}
\centering
\small
\begin{tabular}{p{0.22\textwidth}p{0.68\textwidth}}
\toprule
Plan item & Verified result \\
\midrule
$\KT=0$ source specificity & Correct transport improves SW over wrong source by $0.0386\pm0.0141\sigstar$ and fine-class consistency by $0.0597\pm0.0204\sigstar$. \\
FiLM generality & Correct transport beats shuffled, mean-source, and relation-only controls on SW and semantic consistency at $\KT=1,5,20$. \\
Linear relation & Linear transport is not worse than the nonlinear MLP and improves SW by 0.0093, 0.0112, and 0.0087 at $\KT=1,5,20$. \\
Basis vs. FiLM & Transported generation quality is unresolved as different.  Basis uses 55,344 conditioning parameters versus 118,784 for FiLM (46.6\%). \\
Full fine-tuning & Transport is significantly safer at $\KT=1,5$; full fine-tuning becomes significantly better at $\KT=20$. \\
25\% training tasks & In the frozen one-seed run, basis improves SW by 0.0181--0.0203 and energy MMD by 0.4668--0.4975 over FiLM at $\KT=1,5,20$. \\
\bottomrule
\end{tabular}
\end{table}

Latency, FLOPs, and peak-memory numbers were not present as matched measurements and are not
inferred from parameter count.  The 25\% task result is limited to one training seed.

\paragraph{P5 conclusion.}
The controlled core survives audit: source identity matters at zero target examples, the
relation is approximately linear, its benefit concentrates in the low-data regime, and it
generalises across conditioning mechanisms.  The additive basis is an efficiency result, not
a demonstrated generation-quality winner over FiLM.

\section{P6: Optional Downstream Demographic Augmentation}

P6 was \textbf{not run---gate not met}.  Its strict trigger required a clearly positive P1
correct-versus-wrong family-semantic result.  P1 instead failed the semantic-instrument
validity rule and showed no resolved correct-versus-wrong advantage on the valid secondary
metrics.  Launching or tuning a downstream classifier would therefore violate the plan.

\section{Discussion}

The Fitzpatrick pattern is consistent across four complementary analyses.  First, family
resolution does not separate correct from wrong source coordinates.  Second, wrong diseases
within and outside a family are indistinguishable.  Third, a frozen source-only separability
score does not predict the source advantage.  Fourth, increasing the phenotype gap to V--VI
improves neither task-source specificity nor semantic-instrument validity.  Training
diagnostics independently show that a shared mean coordinate captures nearly all coordinate
gain: shuffled task coordinates and correct coordinates yield effectively identical losses.

This does not imply that demographic transport is impossible, nor does it prove that all
source tasks are equivalent.  The semantic evaluator fails on several held-out families and
diseases, the number of independent tasks is small, and patient identifiers are unavailable.
It does show that the frozen method and available measurements do not justify a real-data
source-specific claim.  In particular, a significant improvement over population-only is
insufficient: wrong-task transport achieves the same improvement, identifying a shared
population offset rather than a task-conditioned relation.

\section{Conclusion}

The executable plan is complete through its registered stopping conditions.  On CIFAR,
\method{} supports the sequence
\[
\boxed{\begin{gathered}
\text{source identity matters}
\;\rightarrow\;
\text{an approximately linear relation exists}
\\[-1pt]
\rightarrow\;
\text{benefit is largest under target scarcity}
\;\rightarrow\;
\text{the principle survives FiLM}.
\end{gathered}}
\]
On Fitzpatrick17k, the corresponding source-specific demographic extension is not supported.
Correct and wrong task sources remain unresolved at disease and family resolution, including
under the stronger V--VI target.  The strongest defensible interpretation is therefore that
the frozen model learns a shared phenotype shift but does not demonstrate disease-specific
coordinate transport.  The invalid semantic instrument requires the headline family result
to remain inconclusive, not equivalent.  No demographic architecture rescue or downstream
augmentation is warranted under the locked protocol.

\appendix
\section{Frozen Configuration and Artifact Map}

\begin{table}[h]
\caption{Frozen family-level training and evaluation configuration.}
\label{tab:config}
\centering
\small
\begin{tabular}{ll}
\toprule
Component & Value \\
\midrule
Image resolution / backbone width & $32\times32$ / 32 \\
Conditioning / coordinate dimension & additive basis / 16 \\
Transport & linear, approximately 0.3K parameters \\
Training & 20,000 steps, AdamW, lr $2\times10^{-4}$ \\
EMA / training seed & 0.999 / 12,345 \\
Source support & $\MS=32$ \\
Target budgets & $\KT=0,1,2,5,10,20$ \\
Generation & DDIM-50, $\eta=0$, 96 images/cell \\
Paired repeats / evaluation seed & 8 / 4,321 \\
Headline target-time refinement & disabled \\
\bottomrule
\end{tabular}
\end{table}

Machine-readable raw cells, family/disease-level values, checkpoint-selection records,
manifests, hashes, commands, and complete tables are stored under
\texttt{experiments/meta\_diffusion\_final/}.  The principal files are
\texttt{protocol\_lock.md}, \texttt{prerequisite\_audit.md},
\texttt{metrics/p1\_family\_aggregate.json},
\texttt{metrics/p2\_hierarchical.json},
\texttt{metrics/p3\_task\_separability.json},
\texttt{metrics/p4\_v\_vi\_aggregate.json}, and
\texttt{runs/execution\_log.md}.  Debug runs are explicitly named and excluded from every
reported estimate.

\end{document}

